\bookchapter{الفصل 7: تحليل شامل للوظائف المتقدمة}

\chapterintro{تتطور التكنولوجيا يوماً بعد يوم، وتتطور وظائف الهاتف باستمرار. اتصال الأقمار الصناعية، والاتصال دون شبكة، وكشف الحوادث، وتقنية SparkLink... هذه الوظائف التي تبدو خيال علمي دخلت بالفعل حياتنا. لنتعرف معاً على هذه التقنيات المتقدمة.}

\booksection{أولاً: اتصال الأقمار الصناعية: الاستغاثة حتى دون إشارة}

اتصال الأقمار الصناعية إنجاز كبير في مجال الهواتف في السنوات الأخيرة. عندما تكون في البرية أو الجبال العميقة، ولا يحتوي هاتفك على أي إشارة، قد يكون اتصال الأقمار الصناعية هو طوق النجاة.

\booksubsection{مبدأ العمل:}

يتصل الهاتف مباشرة بأقمار الاتصالات على المدار لإرسال رسائل الاستغاثة. يدعم حالياً بشكل رئيسي إرسال الرسائل النصية ومعلومات الموقع، ولا يمكن إجراء المكالمات.

\booksubsection{الموديلات المدعومة:}

\begin{itemize}
  \item \textbf{Apple}: سلسلة iPhone 14 وما بعدها تدعم إرسال SOS طارئ عبر الأقمار الصناعية. سلسلة iPhone 17 تدعم بشكل إضافي إرسال iMessage والرسائل النصية عبر الأقمار الصناعية، مع توسيع التغطية لمناطق أكثر.
  \item \textbf{Huawei}: بدأت سلسلة Mate 50 في دعم رسائل الأقمار الصناعية BeiDou (نص فقط)، وتدعم سلسلة Mate 60 Pro بشكل إضافي مكالمات الأقمار الصناعية Tiantong (يمكن إجراء المكالمات)، وتعزّزت سلسلة Mate 70 من قدرات اتصال الأقمار الصناعية بشكل أكبر.
  \item \textbf{Xiaomi/OPPO/vivo}: اعتباراً من 2024، تدعم الموديلات الرائدة تدريجياً وظيفة رسائل الأقمار الصناعية BeiDou.
  \item \textbf{اتجاه التطور}: ينتشر اتصال الأقمار الصناعية من وظيفة حصرية للرائدات المتقدمة إلى الفئات الأدنى، ومن المتوقع أن تدعم الموديلات المتوسطة المدى هذه الوظيفة مستقبلاً. نظام الأقمار الصناعية الصيني BeiDou يحظى بأوسع دعم في الهواتف المحلية.
\end{itemize}

\booksubsection{سيناريوهات الاستخدام:}

\begin{itemize}
  \item المغامرات في الهواء الطلق، والحوادث أثناء تسلق الجبال.
  \item الحوادث في البحر.
  \item انقطاع الاتصالات الأرضية بسبب الكوارث الطبيعية.
\end{itemize}

\booksubsection{طريقة الاستخدام:}

\begin{enumerate}
  \item تأكد من أنك في منطقة مفتوحة، ويمكنك رؤية السماء.
  \item افتح وظيفة اتصال الأقمار الصناعية في الهاتف (عادة في واجهة المكالمات الطارئة).
  \item اتبع إرشادات الشاشة، ووجّه الهاتف نحو القمر الصناعي.
  \item أرسل رسالة الاستغاثة والموقع المُعدّة مسبقاً.
\end{enumerate}

\begin{tcbinfo}{معلومة صغيرة}
يتطلب اتصال الأقمار الصناعية هوائيّاً خاصاً للجهاز، ولا تدعمه جميع الهواتف. كما أن اتصال الأقمار الصناعية عادة ما يتطلب اشتراكاً مدفوعاً، لكن وظيفة الاستغاثة الطارئة تكون مجانية عموماً.

\end{tcbinfo}

\booksection{ثانياً: الاتصال دون شبكة: طريقة جديدة للاتصال قصير المدى}

“الاتصال دون شبكة”يشير إلى تقنية الاتصال المباشر بين هاتفين دون وجود شبكة هاتف محمول أو Wi-Fi.

\booksubsection{المبدأ التقني:}

استخدام بلوتوث أو ترددات راديو مخصصة، لإقامة اتصال نظير إلى نظير ضمن مدى قصير (عادة من بضع مئات الأمتار إلى بضعة كيلومترات).

\booksubsection{سيناريوهات التطبيق:}

\begin{itemize}
  \item \textbf{الأنشطة الخارجية}: التواصل بين أعضاء الفريق أثناء المشي والتخييم.
  \item \textbf{الإنقاذ في الكوارث}: عند انقطاع الاتصالات بعد الزلازل والفيضانات.
  \item \textbf{الفعاليات الكبيرة}: عند ازدحام الناس وضعف الإشارة في الحفلات الموسيقية والفعاليات الرياضية.
\end{itemize}

\booksubsection{المنصات المدعومة:}

\begin{itemize}
  \item \textbf{Apple}: وظائف"NameDrop" و"العثور على الأصدقاء دون اتصال"في iOS 17.
  \item \textbf{Huawei}: تطبيق Changjiang يدعم المكالمات دون شبكة.
  \item \textbf{Xiaomi}: بعض وظائف مركز Xiaomi Miao Xiang.
\end{itemize}

\booksection{ثالثاً: كشف الحوادث: إنقاذ الأرواح في اللحظات الحرجة}

كشف الحوادث وظيفة تستخدم مستشعرات الهاتف للكشف عن الحوادث المرورية الخطيرة، والاتصال تلقائياً برقم الطوارئ.

\booksubsection{مبدأ العمل:}

يستخدم الهاتف عدة مستشعرات مثل مقياس التسارع والجيروسكوب والبارومتر والميكروفون، للحكم بشكل شامل على ما إذا كان قد وقع تصادم خطير. عند الكشف عن حادث، سيقوم الهاتف بـ:

\begin{enumerate}
  \item عرض واجهة عد تنازلي، والسؤال عما إذا كنت بخير.
  \item إذا لم تستجب عند انتهاء العد التنازلي، الاتصال تلقائياً برقم الطوارئ.
  \item إرسال موقعك إلى جهات الاتصال للطوارئ.
\end{enumerate}

\booksubsection{الموديلات المدعومة:}

\begin{itemize}
  \item \textbf{Apple}: iPhone 14 وما بعده.
  \item \textbf{أندرويد}: بعض الموديلات المتقدمة تدعمها تدريجياً.
\end{itemize}

\booksubsection{ملاحظات:}

\begin{itemize}
  \item الوظيفة تحتاج إلى تفعيل لاستخدامها.
  \item قد تُطلق بالخطأ (مثل في الألعاب أو الرياضة العنيفة)، ألغِها في الوقت المناسب.
  \item تعتمد على المستشعرات والخوارزميات، وليست موثوقة 100\%، لكنها طبقة إضافية من الحماية.
\end{itemize}

\booksection{رابعاً: تقنية SparkLink: الاتصال اللاسلكي للجيل القادم}

SparkLink تقنية اتصال لاسلكي من الجيل الجديد قادتها Huawei، وهدفها استبدال بلوتوث وWi-Fi في بعض السيناريوهات. تضاهي أو حتى تتجاوز الحلول الحالية في المؤشرات الرئيسية: تصل سرعة الواجهة الهوائية إلى 900Mbps (أعلى بكثير من 50Mbps لبلوتوث)، وزمن استجابة منخفض حتى 20 $\mu$ s (بلوتوث حوالي 15ms)، وتدعم الشبكة الواحدة 4096 جهازاً متصلاً في الوقت نفسه (بلوتوث عادة أقل من 8).

\booksubsection{الميزات التقنية:}

\begin{itemize}
  \item \textbf{زمن استجابة فائق الانخفاض}: زمن استجابة طرف إلى طرف 20 ميكروثانية فقط (0.02ms)، أقل بمرتبتين من بلوتوث، تزامن الألعاب والموسيقى شبه غير محسوس.
  \item \textbf{سرعة عالية}: سرعة قصوى لرابط واحد 900Mbps، عشرات أضعاف صوت بلوتوث، يمكنها نقل صوت غير مضغوط وحتى فيديو عالي الدقة.
  \item \textbf{استهلاك طاقة منخفض}: استهلاك الطاقة في وضع الاستعداد مماثل لبلوتوث، لكن كفاءة النقل أعلى، واستهلاك الطاقة لكل وحدة بيانات أقل.
  \item \textbf{تغطية واسعة+تزامن عالٍ}: مسافة الاتصال أطول بنسبة 50\% من بلوتوث، وتدعم الشبكة الواحدة أكثر من ألف جهاز متصل في الوقت نفسه.
\end{itemize}

\booksubsection{سيناريوهات التطبيق:}

\begin{itemize}
  \item \textbf{سماعات الرأس اللاسلكية}: زمن استجابة أقل، تزامن الصوت والصورة في الألعاب.
  \item \textbf{الساعات الذكية}: اتصال أكثر استقراراً.
  \item \textbf{المنزل الذكي}: اتصال أجهزة أكثر في الوقت نفسه.
  \item \textbf{التوصيل بالسيارة}: اتصال أكثر سلاسة بين الهاتف ونظام السيارة.
\end{itemize}

\booksubsection{الأجهزة المدعومة:}

حالياً بشكل رئيسي أجهزة منظومة Huawei، مثل سلسلة Huawei Mate 60 وبعض سماعات وساعات Huawei. من المتوقع أن تدعم المزيد من العلامات مستقبلاً.

\textbf{النظرة المستقبلية للتقنية}\par

تمثل تقنية SparkLink اتجاه المستقبل للاتصال اللاسلكي. مع انضمام المزيد من الشركات المصنّعة، قد نرى المزيد من الأجهزة التي تدعم SparkLink، للاستمتاع بتجربة لاسلكية أفضل.

\begin{tcbpractice}{تحقق: هل يدعم هاتفك هذه الوظائف}
\begin{enumerate}
  \item ابحث في الإعدادات عن"قمر صناعي"، لترى ما إذا كانت هناك خيارات متعلقة باتصال الأقمار الصناعية.
  \item ابحث عن"كشف الحوادث" أو "SOS الطارئ"، وتحقق من تفعيلها.
  \item إذا كان هاتفك من Huawei، تحقق مما إذا كان يدعم اتصال أجهزة SparkLink.
  \item تعرف على طرق استخدام هذه الوظائف، فقد تكون مفيدة في اللحظات الحرجة.
\end{enumerate}

\end{tcbpractice}
